
\chapter{Expressions\label{chap:Expressions}}

\hypertarget{def-rhsexpression}{}
Expressions calculate \nativevaluesterm{} from other \nativevaluesterm{}.
We will often refer to expressions defined in this chapter as \rhsexpressions{} to distinguish them
from \assignableexpressions, which are defined in \chapref{AssignableExpressions}.

Expressions can have side effects and can raise exceptions.
Therefore, ASL specifies an evaluation order to ensure that the side-effects/exceptions are predictable (see \secref{EvaluationOrder}).

\ExampleDef{Expressions}
\listingref{Expressions}, \listingref{expr1}, and \listingref{expr2} exemplify the different kinds of expressions,
as indicated by respective comments.
\ASLListing{Examples of expressions}{Expressions}{\definitiontests/Expressions.asl}
\ASLListing{Examples of expressions}{expr1}{\syntaxtests/expr1.asl}
\ASLListing{More examples of expressions}{expr2}{\syntaxtests/expr2.asl}

\ChapterOutline
\begin{itemize}
  \item \FormalRelationsRef{Expressions} defines the formal relations for expressions;
  \item \secref{LiteralExpressions} defines \literalexpressionsterm;
  \item \secref{VariableExpressions} defines \variableexpressionsterm;
  \item \secref{BinaryExpressions} defines \binopexpressionsterm;
  \item \secref{UnaryExpressions} defines \unopexpressionsterm;
  \item \secref{ConditionalExpressions} defines \condexpressionsterm;
  \item \secref{CallExpressions} defines \callexpressionsterm;
  \item \secref{SlicingExpressions} defines \slicingexpressionsterm;
  \item \secref{ArrayAccessExpressions} defines \arrayaccessexpressionsterm;
  \item \secref{FieldReadingExpressions} defines field reading expressions;
  \item \secref{MultiFieldReadingExpressions} defines multi-field reading expressions;
  \item \secref{AssertingTypeConversionExpressions} defines asserting type conversion expressions;
  \item \secref{PatternMatchingExpressions} defines pattern matching expressions;
  \item \secref{ArbitraryValueExpressions} defines arbitrary value expressions;
  \item \secref{StructuredTypeConstructionExpressions} defines structured type construction expressions;
  \item \secref{TupleExpressions} defines tuple expressions;
  \item \secref{ParenthesizedExpressions} defines parenthesized expressions;
  \item \secref{ArrayConstructionExpressions} defines array construction expressions;
  \item \secref{SideEffectFreeExpressions} defines \sideeffectfreeexpressionsterm{};
  \item \secref{ExprList} defines the dynamic semantics of a list of expressions.
\end{itemize}

%%%%%%%%%%%%%%%%%%%%%%%%%%%%%%%%%%%%%%%%%%%%%%%%%%%%%%%%%%%%%%%%%%%%%%%%%%%%%%%%
\FormalRelationsDef{Expressions}
\paragraph{Syntax:} Expressions are grammatically derived from $\Nexpr$.

\paragraph{Abstract Syntax:} Expressions are derived in the abstract syntax from $\expr$,
and generated by $\buildexpr$.
\NotImportedToASLSpecYet{
\hypertarget{build-expr}{}
The function
\[
  \buildexpr(\overname{\parsenode{\Nexpr}}{\vparsednode}) \;\aslto\; \overname{\expr}{\vastnode}
  \cup \overname{\TBuildError}{\BuildErrorConfig}
\]
transforms an expression parse node $\vparsednode$ into an expression AST node $\vastnode$.
\ProseOtherwiseBuildError
} % END_OF_UNIMPORTED_RELATION


\paragraph{Typing:}
All expressions have a unique type (which can be a \tupletypeterm{}).
\RenderRelation{annotate_expr}
The annotation rewrites the input expression in the following case, making the annotation of statements simpler:
variables with constant values are substituted by their constant values.

\paragraph{Semantics:}
\RenderRelation{eval_expr}
\hypertarget{def-literalexpressionterm}{}
\section{Literal Expressions\label{sec:LiteralExpressions}}
A literal expression represents a literal as an expression.

\ASLListing{Literal Expressions}{literalssemantics}{\semanticstests/SemanticsRule.Lit.asl}

\subsection{Syntax}
\begin{flalign*}
\Nexpr \derives\  & \Nvalue &
\end{flalign*}

\subsection{Abstract Syntax}
\RenderTypes[remove_hypertargets]{expr_literal}


\ASTRuleDef{ELit}
\begin{mathpar}
\inferrule{}{
  \buildexpr(\overname{\Nexpr(\punnode{\Nvalue})}{\vparsednode}) \astarrow
  \overname{\ELiteral(\astof{\vvalue})}{\vastnode}
}
\end{mathpar}

\subsection{Typing}
\TypingRuleDef{ELit}
\ExampleDef{Typing of Literal Expressions}
In \listingref{literalssemantics}, each of the expressions \texttt{3}
is annotated with the type \verb|integer{3}|.

\RenderProseAndFormally{annotate_expr_ELit}
\CodeSubsection{\ELitBegin}{\ELitEnd}{../Typing.ml}

\subsection{Semantics}
\SemanticsRuleDef{ELit}
\ExampleDef{Evaluation of Literal Expressions}
In \listingref{literalssemantics}, each of the expressions \texttt{3} evaluates to the \nativevalueterm{} $\nvint(3)$.

\RenderProseAndFormally{eval_expr_ELit}
\CodeSubsection{\EvalELitBegin}{\EvalELitEnd}{../Interpreter.ml}

\hypertarget{def-variableexpressionterm}{}
\section{Variable Expressions\label{sec:VariableExpressions}}
A variable expression consists of an identifier for a storage element.

\subsection{Syntax}
\begin{flalign*}
\Nexpr \derives\ & \Tidentifier &
\end{flalign*}

\subsection{Abstract Syntax}
\RenderTypes[remove_hypertargets]{expr_var}


\ASTRuleDef{EVAR}
\begin{mathpar}
  \inferrule{}{
  \buildexpr(\overname{\Nexpr(\Tidentifier(\id))}{\vparsednode}) \astarrow
  \overname{\EVar(\id)}{\vastnode}
}
\end{mathpar}

\subsection{Typing}
\TypingRuleDef{EVar}
\ExampleDef{Typing Variable Expressions}
All of the variable expressions in \listingref{expressions-evar}
are well-typed.
\ASLListing{Variable expressions}{expressions-evar}{\typingtests/TypingRule.EVar.asl}

The variable \texttt{t} in \listingref{expressions-evar-undefined} is undefined.
\ASLListing{An undefined variable}{expressions-evar-undefined}{\typingtests/TypingRule.EVar.undefined.asl}

\RenderProseAndFormally{annotate_expr_EVar}

\subsection{Semantics}
\SemanticsRuleDef{EVar}

\ExampleDef{Evaluation of a Local Variable Expression}
In \listingref{localvarsemantics}, the evaluation of \texttt{x} within \texttt{assert x == 3;}
applies the rule \SemanticsRuleRef{EVar}.LOCAL.
\ASLListing{Semantics of local variables}{localvarsemantics}{\semanticstests/SemanticsRule.ELocalVar.asl}

\ExampleDef{Evaluation of a Global Variable Expression}
In \listingref{globalvarsemantics}, the evaluation of~\texttt{global\_x} within~\texttt{assert global\_x == 3;}
applies the rule \SemanticsRuleRef{EVar}.GLOBAL.
\ASLListing{Semantics of global variables}{globalvarsemantics}{\semanticstests/SemanticsRule.EGlobalVar.asl}

\RenderProseAndFormally{eval_expr_EVar}
\CodeSubsection{\EvalEVarBegin}{\EvalEVarEnd}{../Interpreter.ml}

\subsubsection{Comments}
When a global variable, $\vx$, exists, the type system
forbids having $\vx$ as a local variable.
This is enforced by \TypingRuleRef{LDVar} in the Chapter ``Typing of Local Declarations'',
and
\TypingRuleRef{DeclareGlobalStorage} and \TypingRuleRef{DeclareOneFunc},
both in the Chapter ``Typing of Global Declarations''.

\hypertarget{def-binopexpressionterm}{}
\section{Binary Expressions\label{sec:BinaryExpressions}}
Binary expressions apply a binary operator to two sub-expressions.

\RequirementDef{OperatorPrecedence}
Operator precedence is used to disambiguate binary expressions.

Given two binary operators $\opone$ and $\optwo$, an expression of the form
$x\ \opone\ y\ \optwo\ z$, is interpreted as
$(x\ \opone\ y)\ \optwo\ z$
if $\opone$ has higher precedence than $\optwo$ and as $x\ \opone\ (y\ \optwo\ z)$ if $\opone$ has lower precedence than $\optwo$.
If $\opone$ is associative and $\opone$ = $\optwo$ then the expression may be interpreted as either
$(x\ \opone\ y)\ \optwo\ z$ or as $x\ \opone\ (y\ \optwo\ z)$ since there is no difference.
Otherwise it is a static operator precedence error (\BinopPrecedence).
See \ASTRuleRef{EBinop} for the formal details.

The operator classes and precedence order are defined by the following table
(a subset of the table in \secref{PriorityAndAssociativity},
which contains tokens that are not operators, but are needed for parsing disambiguation):

\begin{center}
\begin{tabular}{lll}
\hline
\textbf{Precedence} & \textbf{Class} & \textbf{Operators}\\
\hline
7 (Highest) & \textbf{Membership} & \Tin\\
6 & \textbf{Unary}                & \Tminus, \Tnot, \Tbnot\\
5 & \textbf{Power}                & \Tpow\\
4 & \textbf{Mul--Div--Shift}      & \Tmul, \Tdiv, \Tdivrm, \Trdiv, \Tmod, \Tshl, \Tshr\\
3 & \textbf{Add--Sub--Logic}      & \Tplus, \Tminus, \Tor, \Tand, \Txor, \Tcoloncolon, \Tplusplus\\
2 & \textbf{Relational}           & \Tgt, \Tgeq, \Tlt, \Tleq\\
1 & \textbf{(In-)Equality}        & \Teqop, \Tneq\\
0 & \textbf{Boolean}              & \Tbor, \Tband, \Tbimpl, \Tbeq\\
\hline
\end{tabular}
\end{center}

The set of \emph{associative binary operators} consists of the following:
$\ADD$,
$\MUL$,
$\BAND$,
$\BOR$,
$\AND$,
$\OR$,
$\XOR$,
$\BEQ$,
$\BVCONCAT$,
$\STRCONCAT$.

The specification in \listingref{BinaryExpressions} shows examples of binary expressions
and the effect of operator precedence.
\ASLListing{Binary expressions and operator precedence}{BinaryExpressions}{\definitiontests/Guide.OperatorPrecedence.asl}

\subsection{Syntax}
\begin{flalign*}
\Nexpr \derives\  & \Nexpr \parsesep \Nbinop \parsesep \Nexpr &
\end{flalign*}

\subsection{Abstract Syntax}
\RenderTypes[remove_hypertargets]{expr_binop}

\ASTRuleDef{EBinop}
The following rule constructs a binary expression AST from a parse tree,
which possibly contains parentheses,
ensuring that \RequirementRef{OperatorPrecedence} holds.

The specification in \listingref{BinaryExpressions}
shows examples where \RequirementRef{OperatorPrecedence} holds.
%
The specifications in
\listingref{EBinop-bad1},
\listingref{EBinop-bad2},
\listingref{EBinop-bad3},
\listingref{EBinop-bad4}, and
\listingref{EBinop-bad5}
show examples that violate \RequirementRef{OperatorPrecedence},
which results in a \builderrorterm.
\ASLListing{Violation of operator precedence}{EBinop-bad1}{\syntaxtests/ASTRule.EBinop.bad1.asl}
\ASLListing{Violation of operator precedence}{EBinop-bad2}{\syntaxtests/ASTRule.EBinop.bad2.asl}
\ASLListing{Violation of operator precedence}{EBinop-bad3}{\syntaxtests/ASTRule.EBinop.bad3.asl}
\ASLListing{Violation of operator precedence}{EBinop-bad4}{\syntaxtests/ASTRule.EBinop.bad4.asl}
\ASLListing{Violation of operator precedence}{EBinop-bad5}{\syntaxtests/ASTRule.EBinop.bad5.asl}

\begin{mathpar}
  \inferrule{
    \buildexpr(\veone) \astarrow \astversion{\veone} \OrBuildError\hva\\\\
    \buildexpr(\vetwo) \astarrow \astversion{\vetwo} \OrBuildError\hva\\\\
    \checknotsameprec(\astof{\vbinop}, \astversion{\veone}) \astarrow \True \OrBuildError\hva\\\\
    \checknotsameprec(\astof{\vbinop}, \astversion{\vetwo}) \astarrow \True \OrBuildError
  }{
    {
      \begin{array}{r}
  \buildexpr(\overname{\Nexpr(\namednode{\veone}{\Nexpr}, \punnode{\Nbinop}, \namednode{\vetwo}{\Nexpr})}{\vparsednode}) \astarrow\\
  \overname{\EBinop(\astversion{\veone}, \astof{\vbinop}, \astversion{\vetwo})}{\vastnode}
      \end{array}
    }
}
\end{mathpar}

\ASTRuleDef{CheckNotSamePrec}
\hypertarget{build-binopprec}{}
We define the helper function
\[
  \binopprec(\overname{\binop}{\op}) \aslto \N
\]
which assigns a precedence level to each binary operator, $\op$, as defined below:
\begin{mathpar}
\inferrule{}{
  {
  \binopprec(\op) \astarrow
  \begin{cases}
    5 & \text{if }\op = \POW\\
    4 & \text{if }\op \in \{\MUL, \DIV, \DIVRM, \RDIV, \MOD, \SHL, \SHR\}\\
    3 & \text{if }\op \in \left\{
        \begin{array}{l}
          \ADD, \SUB, \OR, \XOR, \\
          \AND, \BVCONCAT, \STRCONCAT
        \end{array}
        \right\}\\
    2 & \text{if }\op \in \{\GT, \GE, \LT, \LE\}\\
    1 & \text{if }\op \in \{\EQ, \NE\}\\
    0 & \text{if }\op \in \{\BOR, \BAND, \IMPL, \BEQ \}\\
  \end{cases}
  }
}
\end{mathpar}

\NotImportedToASLSpecYet{
\hypertarget{build-checknotsameprec}{}
The function
\[
\checknotsameprec(\overname{\binop}{\op} \aslsep \overname{\expr}{\ve})
\aslto \{\True\} \cup \overname{\TBuildError}{\BuildErrorConfig}
\]
checks whether the expression AST node $\ve$ is a binary operator, either with a different operator to $\op$ but the same \emph{precedence}, or the same non-associative operator $\op$.
In either case, it is considered an error.
Surrounding $\ve$ by parenthesis fixes the error.
} % END_OF_UNIMPORTED_RELATION


For example, \texttt{a + b + c} is considered legal, since the same binary operator (\texttt{+})
is used, whereas \texttt{a - b - c} and \texttt{a + b - c} are considered illegal.
For the first, this is because $\SUB$ is not associative.
For the second, this is because $\ADD$ and $\SUB$ have the
same precedence ($3$). To fix these, we can surround one of the subexpressions with parenthesis,
for example: \texttt{(a - b) - c} and \texttt{(a + b) - c}.

\ProseParagraph
\OneApplies
\begin{itemize}
  \item \AllApplyCase{not\_binop}
  \begin{itemize}
    \item $\ve$ is not a binary operation expression;
    \item the result is $\True$.
  \end{itemize}

  \item \AllApplyCase{binop\_same}
  \begin{itemize}
    \item $\ve$ is a binary operation expression for the operator $\opp$;
    \item $\opp$ is the same operator as $\op$;
    \item checking whether $\op$ is an associative operator yields $\True$\ProseTerminateAs{\BinopPrecedence}.
  \end{itemize}

  \item \AllApplyCase{binop\_different}
  \begin{itemize}
    \item $\ve$ is a binary operation expression for the operator $\opp$;
    \item $\opp$ is not the same operator as $\op$;
    \item checking whether $\op$ and $\opp$ have different precedence levels yields $\True$\ProseTerminateAs{\BinopPrecedence}.
  \end{itemize}
\end{itemize}

\FormallyParagraph
\begin{mathpar}
\inferrule[not\_binop]{
  \astlabel(\ve) \neq \EBinop
}{
  \checknotsameprec(\op, \ve) \astarrow \True
}
\end{mathpar}

\begin{mathpar}
\inferrule[binop\_same]{
  \op = \opp \\
  {
  \vassociativeoperators \eqdef \left\{
    \begin{array}{l}
    \ADD, \\
    \MUL, \\
    \BAND, \\
    \BOR, \\
    \AND, \\
    \OR, \\
    \XOR, \\
    \BEQ, \\
    \BVCONCAT, \\
    \STRCONCAT
    \end{array}
    \right\}
  } \\
  {
  \techeck(\op \in \vassociativeoperators, \BinopPrecedence) \typearrow \True \OrBuildError
  }
}{
  \checknotsameprec(\op, \overname{\EBinop(\opp, \Ignore, \Ignore)}{\ve}) \astarrow \True
}
\end{mathpar}

\begin{mathpar}
\inferrule[binop\_different]{
  \op \neq \opp \\
  \becheck(\binopprec(\op) \neq \binopprec(\opp), \BinopPrecedence) \astarrow \True \OrBuildError
}{
  \checknotsameprec(\op, \overname{\EBinop(\opp, \Ignore, \Ignore)}{\ve}) \astarrow \True
}
\end{mathpar}

\subsection{Typing}
\TypingRuleDef{EBinop}
\ExampleDef{An ill-typed Binary Expression}
The specification in \listingref{typing-binoperror} is ill-typed, since,
according to the signature of $\divint$, \Tdiv{} can only be applied to \integertypesterm{}
whereas \verb|0.0| is a \realtypeterm.
\ASLListing{An ill-typed binary expression}{typing-binoperror}{\typingtests/TypingRule.EBinop.bad.asl}

\ExampleDef{Comparing Well-constrained Integers}
The specification in \listingref{EBinop} is well-typed.
Specifically, it is legal to compare \wellconstrainedintegertypesterm{}
with different lists of constraints and a \wellconstrainedintegertypeterm{} to an \unconstrainedintegertypeterm{}.
\ASLListing{Comparing well-constrained integers}{EBinop}{\typingtests/TypingRule.EBinop.asl}

\ExampleDef{Typing Bitvector Concatenation}
The specification in \listingref{EBinop-2} is well-typed, since
the sum of widths of the two concatenated \bitvectortypesterm{}
is equal to the width of the \bitvectortypeterm{} of the return type.
\ASLListing{Concatenating bitvectors}{EBinop-2}{\typingtests/TypingRule.EBinop2.asl}

\RenderProseAndFormally{annotate_expr_EBinop}
\CodeSubsection{\BinopBegin}{\BinopEnd}{../Typing.ml}

\subsection{Semantics}
\SemanticsRuleDef{BinopAnd}
\ExampleDef{Evaluation of Boolean And Expressions}
In \listingref{andsemantics},
the expression \texttt{FALSE \&\& fail()} evaluates to the value \texttt{FALSE}. Notice that the function \texttt{fail} is never called.
\ASLListing{Semantics of conjunction}{andsemantics}{\semanticstests/SemanticsRule.EBinopAndFalse.asl}

\RenderProseAndFormally{eval_expr_BinopAnd}
\CodeSubsection{\EvalBinopAndBegin}{\EvalBinopAndEnd}{../Interpreter.ml}

\subsubsection{Comments}
The evaluation via the rule above ensures that $\veone$ is evaluated first and
only if it evaluates to $\True$ is $\vetwo$ evaluated.

 % TODO: add table


Conditional expressions and the operations \texttt{\&\&}, \texttt{||},
\texttt{==>} provide a short-circuit evaluation mechanism:

% This is related to the note under
\begin{itemize}
\item the first operand of \texttt{if} is always evaluated but only one of the
remaining operands is evaluated;
\item if the first operand of \texttt{and\_bool} is $\False$, then the second operand is not evaluated;
\item if the first operand of \texttt{or\_bool} is $\True$, then the second operand is not evaluated; and,
\item if the first operand of \texttt{implies\_bool} is $\False$, then the
second operand is not evaluated.
\end{itemize}

However, note that relying on this short-circuit evaluation can be confusing
for readers of ASL specifications and as a consequence it is recommended that
an if-statement is used to achieve the same effect.

\SemanticsRuleDef{BinopOr}
\ExampleDef{Evaluation of Boolean Or Expressions}
In \listingref{semantics-binopor}, the expression \texttt{(0 == 1) || (1 == 1)} evaluates to the value \True.
\ASLListing{Evaluating a disjunction expression}{semantics-binopor}{\semanticstests/SemanticsRule.EBinopOrTrue.asl}

\RenderProseAndFormally{eval_expr_BinopOr}
\CodeSubsection{\EvalBinopOrBegin}{\EvalBinopOrEnd}{../Interpreter.ml}

The evaluation via the rule above ensures that $\veone$ is evaluated first and only if
it evaluates to $\False$, is $\vetwo$ evaluated.

\subsubsection{Comments}
 % TODO: add table


Conditional expressions and the operations \texttt{\&\&}, \texttt{||},
\texttt{==>} provide a short-circuit evaluation mechanism:

% This is related to the note under
\begin{itemize}
\item the first operand of \texttt{if} is always evaluated but only one of the
remaining operands is evaluated;
\item if the first operand of \texttt{and\_bool} is $\False$, then the second operand is not evaluated;
\item if the first operand of \texttt{or\_bool} is $\True$, then the second operand is not evaluated; and,
\item if the first operand of \texttt{implies\_bool} is $\False$, then the
second operand is not evaluated.
\end{itemize}

However, note that relying on this short-circuit evaluation can be confusing
for readers of ASL specifications and as a consequence it is recommended that
an if-statement is used to achieve the same effect.

\SemanticsRuleDef{BinopImpl}
\ExampleDef{Evaluation of Implication Expressions}
In \listingref{semantics-binopimpl},
the expression \texttt{(0 == 1) ==> (1 == 0)} evaluates to the value \True, according to the definition of implication.
\ASLListing{Evaluating an implication expression}{semantics-binopimpl}{\semanticstests/SemanticsRule.EBinopImplExFalso.asl}

\RenderProseAndFormally{eval_expr_BinopImpl}
\CodeSubsection{\EvalBinopImplBegin}{\EvalBinopImplEnd}{../Interpreter.ml}
The evaluation via the rule above ensures that $\veone$ is evaluated first and only if
it evaluates to \True, is $\vetwo$ evaluated.


Conditional expressions and the operations \texttt{\&\&}, \texttt{||},
\texttt{==>} provide a short-circuit evaluation mechanism:

\begin{itemize}
\item the first operand of \texttt{if} is always evaluated but only one of the
remaining operands is evaluated;
\item if the first operand of \texttt{and\_bool} is $\False$, then the second operand is not evaluated;
\item if the first operand of \texttt{or\_bool} is $\True$, then the second operand is not evaluated; and,
\item if the first operand of \texttt{implies\_bool} is $\False$, then the
second operand is not evaluated.
\end{itemize}

However, note that relying on this short-circuit evaluation can be confusing
for readers of ASL specifications and as a consequence it is recommended that
an if-statement is used to achieve the same effect.

\SemanticsRuleDef{Binop}
\ExampleDef{Evaluation of Binary Expressions}
In \listingref{semantics-binopassert},
the expression \texttt{3 + 2} evaluates to the value \texttt{5}.
\ASLListing{Evaluating a binary expression}{semantics-binopassert}{\semanticstests/SemanticsRule.EBinopPlusAssert.asl}

\RenderProseAndFormally{eval_expr_Binop}
\CodeSubsection{\EvalBinopBegin}{\EvalBinopEnd}{../Interpreter.ml}

The rule above applies to many binary operators, including $\EQ$ (which is used for \texttt{<=>}
as well as \texttt{==}).

\hypertarget{def-unopexpressionterm}{}
\section{Unary Expressions\label{sec:UnaryExpressions}}
\subsection{Syntax}
\begin{flalign*}
\Nexpr \derives\  & \Nunop \parsesep \Nexpr &
\end{flalign*}

\subsection{Abstract Syntax}
\RenderTypes[remove_hypertargets]{expr_unop}


\ASTRuleDef{EUnop}
\begin{mathpar}
  \inferrule{
    \buildexpr(\vexpr) \astarrow \astversion{\vexpr} \OrBuildError
  }{
  \buildexpr(\overname{\Nexpr(\punnode{\Nunop}, \vexpr : \Nexpr)}{\vparsednode}) \astarrow
  \overname{\EUnop(\astof{\vunop}, \astversion{\vexpr})}{\vastnode}
}
\end{mathpar}

\subsection{Typing}
\TypingRuleDef{EUnop}
\ExampleRef{Applying Unary Operations to Types} shows examples
of well-typed unary operation expressions.

\RenderProseAndFormally{annotate_expr_EUnop}
\CodeSubsection{\UnopBegin}{\UnopEnd}{../Typing.ml}

\subsection{Semantics}
\SemanticsRuleDef{Unop}
\ExampleDef{Evaluation of a Unary Operation Expression}
In \listingref{semantics-unopassert},
the expression \texttt{NOT '1010'} evaluates to the value \texttt{'0101'}.
\ASLListing{Evaluating a unary operation expression}{semantics-unopassert}{\semanticstests/SemanticsRule.EUnopAssert.asl}

\RenderProseAndFormally{eval_expr_Unop}
\CodeSubsection{\EvalUnopBegin}{\EvalUnopEnd}{../Interpreter.ml}

\hypertarget{def-conditionexpressionterm}{}
\section{Conditional Expressions\label{sec:ConditionalExpressions}}
\subsection{Syntax}
\begin{flalign*}
\Nexpr \derives\  & \Tif \parsesep \Nexpr \parsesep \Tthen \parsesep \Nexpr \parsesep \Telse \parsesep \Nexpr &\\
\end{flalign*}

\subsection{Abstract Syntax}
\RenderTypes[remove_hypertargets]{expr_cond}


\ASTRuleDef{ECond}
\begin{mathpar}
  \inferrule{
    \buildexpr(\vcondexpr) \astarrow \astversion{\vcondexpr} \OrBuildError\hva\\\\
    \buildexpr(\vthenexpr) \astarrow \astversion{\vthenexpr} \OrBuildError\hva\\\\
    \buildexpr(\velseexpr) \astarrow \astversion{\velseexpr} \OrBuildError\hva\\\\
  }{
    {
      \begin{array}{r}
  \buildexpr\left(\overname{\Nexpr\left(
    \begin{array}{l}
    \Tif, \namednode{\vcondexpr}{\Nexpr}, \Tthen, \\
    \wrappedline\ \namednode{\vthenexpr}{\Nexpr}, \Telse, \namednode{\velseexpr}{\Nexpr}
    \end{array}
    \right)}{\vparsednode}\right) \astarrow\\
  \overname{\ECond(\astversion{\vcondexpr}, \astversion{\vthenexpr}, \astversion{\velseexpr})}{\vastnode}
      \end{array}
    }
}
\end{mathpar}

\subsection{Typing}
\TypingRuleDef{ECond}
\ExampleDef{Typing Conditional Expressions}
The specification in \listingref{ECond} is well-typed and shows
an example of typing a conditional expression.
\ASLListing{Typing conditional expressions}{ECond}{\typingtests/TypingRule.ECond.asl}

\ExampleRef{Lowest Common Ancestor Examples} shows more examples of typing conditional expressions.

\RenderProseAndFormally{annotate_expr_ECond}
\CodeSubsection{\ECondBegin}{\ECondEnd}{../Typing.ml}


\subsection{Semantics}
\SemanticsRuleDef{ECond}
\ExampleDef{Evaluation of Conditional Expressions}
In \listingref{semantics-econdfalse},
the expression \texttt{if FALSE then Return42() else 3} evaluates to the value \texttt{3}.
\ASLListing{Evaluating a conditional expression yielding the result of the \texttt{else} subexpression}{semantics-econdfalse}{\semanticstests/SemanticsRule.ECondFALSE.asl}

\ExampleDef{Evaluation of a Non-deterministic Conditional Expression}
In \listingref{semantics-econdarbitrary},
the expression \texttt{if ARBITRARY: boolean then 3 else Return42()} will
evaluate to either \texttt{3} or \texttt{Return42()}, depending on the
(non-deterministic) result of \texttt{ARBITRARY: boolean}.
\ASLListing{Evaluating a conditional expression with non-deterministic choice}{semantics-econdarbitrary}
{\semanticstests/SemanticsRule.ECondARBITRARY3or42.asl}

% Transliteration note: the code uses an optimized semantics for the case where
% true and false sub-expressions do not contain function calls. Since this is an
% optimization, we do not document it.
\RenderProseAndFormally{eval_expr_ECond}
\CodeSubsection{\EvalECondBegin}{\EvalECondEnd}{../Interpreter.ml}

\hypertarget{def-callexpressionterm}{}
\section{Call Expressions\label{sec:CallExpressions}}
\ASLListing{Call expressions}{semantics-ecall}{\semanticstests/SemanticsRule.ECall.asl}

\subsection{Syntax}
\begin{flalign*}
\Nexpr \derives\  & \Ncall &
\end{flalign*}

\begin{flalign*}
\Ncall \derives \
     & \Tidentifier \parsesep \PlistZero{\Nexpr} &\\
  |\ & \Tidentifier \parsesep \Tlbrace \parsesep \ClistOne{\Nexpr} \parsesep \Trbrace &\\
  |\ & \Tidentifier \parsesep \Tlbrace \parsesep \ClistOne{\Nexpr} \parsesep \Trbrace \parsesep \PlistZero{\Nexpr} &
\end{flalign*}

\subsection{Abstract Syntax}
\RenderTypes[remove_hypertargets]{expr_call}

\ASTRuleDef{Call}
\texthypertarget{build-call}
The function
\[
  \buildcall(\overname{\parsenode{\Ncall}}{\vparsednode}) \astarrow \overname{\call}{\vastnode} \cup \overname{\TBuildError}{\BuildErrorConfig}
\]
constructs a call expression AST from a parse tree for a subprogram call.
\ProseOtherwiseBuildError

\begin{mathpar}
\inferrule{
  \buildplist[\buildexpr](\vargs) \astarrow \vargasts
}{
  \buildcall(\overname{\Ncall(\Tidentifier(\id), \namednode{\vargs}{\PlistZero{\Nexpr}})}{\vparsednode}) \astarrow \\
  { \overname{\left\{
      \begin{array}{rcl}
        \callname &:& \id,\\
        \callparams &:& \emptylist,\\
        \callargs &:& \vargasts,\\
        \callcalltype &:& \STFunction
      \end{array}
    \right\}}{\vastnode} }
}
\hva\and
\inferrule{
  \buildlist[\buildexpr](\vparams) \astarrow \astversion{\vparams} \\
}{
  \buildcall(\overname{\Ncall(\Tidentifier(\id), \Tlbrace, \namednode{\vparams}{\ClistOne{\Nexpr}}, \Trbrace)}{\vparsednode}) \astarrow \\
  { \overname{\left\{
      \begin{array}{rcl}
              \callname &:& \id,\\
              \callparams &:& \astversion{\vparams},\\
              \callargs &:& \emptylist,\\
              \callcalltype &:& \STFunction
      \end{array}
    \right\}}{\vastnode} }
}
\hva\and
\inferrule{
  \buildplist[\buildexpr](\vargs) \astarrow \vargasts \\
  \buildlist[\buildexpr](\vparams) \astarrow \astversion{\vparams} \\
}{
  \buildcall(\overname{\Ncall(\Tidentifier(\id), \Tlbrace, \namednode{\vparams}{\ClistOne{\Nexpr}}, \Trbrace, \namednode{\vargs}{\PlistZero{\Nexpr}})}{\vparsednode}) \astarrow \\
  { \overname{\left\{
      \begin{array}{rcl}
              \callname &:& \id,\\
              \callparams &:& \astversion{\vparams},\\
              \callargs &:& \vargasts,\\
              \callcalltype &:& \STFunction
      \end{array}
    \right\}}{\vastnode} }
}
\end{mathpar}

\ASTRuleDef{ECall}
\begin{mathpar}
\inferrule{
  \buildcall(\vcall) \astarrow \astversion{\vcall} \OrBuildError\hva
}{
  \buildexpr(\overname{\Nexpr(\namednode{\vcall}{\Ncall})}{\vparsednode}) \astarrow
  \overname{\ECall(\astof{\vcall})}{\vastnode}
}
\end{mathpar}

\subsection{Typing}
\TypingRuleDef{ECall}
\ExampleDef{Typing Call Expressions}
\listingref{semantics-ecall} shows call expressions (on the right-hand-side of assignments)
and the inferred types of the expressions, as type annotations on the left-hand-side variables.

\RenderProseAndFormally{annotate_expr_ECall}
\CodeSubsection{\ECallBegin}{\ECallEnd}{../Typing.ml}


\subsection{Semantics}
\SemanticsRuleDef{ECall}
\ExampleDef{Evaluation of Call Expressions}
\listingref{semantics-ecall} shows call expressions (on the right-hand-side of assignments)
and the values they evaluate to, as assertions on the values assigned.

\RenderProseAndFormally{eval_expr_ECall}
\CodeSubsection{\EvalECallBegin}{\EvalECallEnd}{../Interpreter.ml}

\hypertarget{def-slicingexpressionsterm}{}
\section{Slicing Expressions\label{sec:SlicingExpressions}}
This section details the high-level form of the syntax and abstract syntax of slicing expressions,
and defines the semantics of bitvector slices.
The details of the various types of bitvector slices are deferred to \chapref{BitvectorSlicing}.

The well-typed specification in \listingref{Bitvectorslices2} shows examples of bitvector slices.
\ASLListing{Slicing Expressions}{Bitvectorslices2}{\definitiontests/Bitvector_slices2.asl}

\subsection{Syntax}
\begin{flalign*}
\Nexpr \derives\ & \Nexpr \parsesep \Nslices &
\end{flalign*}

\subsection{Abstract Syntax}
\RenderTypes[remove_hypertargets]{expr_slice}


\ASTRuleDef{ESlice}
\begin{mathpar}
\inferrule{
  \buildexpr(\vexpr) \astarrow \astversion{\vexpr} \OrBuildError\hva\\\\
  \buildslices(\vslices) \astarrow \astversion{\vslices} \OrBuildError
}{
  \buildexpr(\overname{\Nexpr(\vexpr: \Nexpr, \vslices: \Nslices)}{\vparsednode}) \astarrow
  \overname{\ESlice(\astversion{\vexpr}, \astversion{\vslices})}{\vastnode}
}
\end{mathpar}

\subsection{Typing}
\TypingRuleDef{ESlice}
\ExampleDef{Typing of Slicing Expressions}
In \listingref{semantics-eslice}, the type inferred for the
slicing expression \verb|'1 1110 000'[6:3, 7]| is
\verb|bits(5){}|.
That is, a bitvector of width \verb|5| without any bitfields.
%
The same type is inferred for the slicing expressions
\verb|240[6:3, 7]| and \verb|496[6:3, 7]|
(\verb|240| is equivalent to \verb|'1 111 0 000'|
and \verb|496| is equivalent to \verb|'11 111 0 000'|).

\ExampleDef{Ill-typed Slicing Expressions}
The specifications in \listingref{ESlice-bad1} and \listingref{ESlice-bad2}
show examples of ill-typed slicing expressions.
\ASLListing{Ill-typed slicing expression due to incorrect width}{ESlice-bad1}{\typingtests/TypingRule.ESlice.bad1.asl}
\ASLListing{Ill-typed slicing expression due to incorrect type}{ESlice-bad2}{\typingtests/TypingRule.ESlice.bad2.asl}

\RenderProseAndFormally{annotate_expr_ESlice}
\CodeSubsection{\ESliceBegin}{\ESliceEnd}{../Typing.ml}
\subsubsection{Comments}
The width of \slices\ might be a symbolic expression if one of the
widths references a \texttt{let} identifier with a non-compile-time-constant
initializer expression.


\subsection{Semantics}
\SemanticsRuleDef{ESlice}
\ExampleDef{Evaluation of Slicing Expressions}
In \listingref{semantics-eslice}, the slicing expression \verb|'1 1110 000'[6:3, 7]|
evaluates to the bitvector value \texttt{'1110 1'},
same as the slicing expressions
\verb|240[6:3, 7]| and \verb|496[6:3, 7]|.

\ASLListing{Slicing Expressions}{semantics-eslice}{\semanticstests/SemanticsRule.ESlice.asl}

Note that in the following definition,
the function $\readfrombitvector$ takes care of converting integers to bitvectors.

\RenderProseAndFormally{eval_expr_ESlice}
\CodeSubsection{\EvalESliceBegin}{\EvalESliceEnd}{../Interpreter.ml}

\hypertarget{def-arrayaccessexpressionsterm}{}
\hypertarget{def-getarrayexpressionterm}{}
\section{Array Access Expressions\label{sec:ArrayAccessExpressions}}
This section details the syntax, abstract syntax, semantics, and typing of array read expressions.

\subsection{Syntax}
\begin{flalign*}
\Nexpr \derives\ & \Nexpr \parsesep \Tllbracket \parsesep \Nexpr \parsesep \Trrbracket &
\end{flalign*}

\subsection{Abstract Syntax}
\RenderTypes[remove_hypertargets]{expr_getarray}

\ASTRuleDef{EGetArray}
\begin{mathpar}
\inferrule{
  \buildexpr(\veone) \astarrow \astversion{\veone} \OrBuildError\hva\\\\
  \buildexpr(\vetwo) \astarrow \astversion{\vetwo} \OrBuildError
}{
  \buildexpr(\overname{\Nexpr(\namednode{\veone}{\Nexpr}, \Tllbracket, \namednode{\vetwo}{\Nexpr}, \Trrbracket)}{\vparsednode}) \astarrow
  \overname{\EGetArray(\astversion{\veone}, \astversion{\vetwo})}{\vastnode}
}
\end{mathpar}

\TypingRuleDef{EGetArray}
\hypertarget{def-arrayaccess}{}
\begin{definition}[Array Access]
We refer to a \rhsexpression{} of the form \verb|b[[i]]|,
where $b, i$ are subexpressions, as an \arrayaccessterm\ expression.
We refer to $b$ and $i$ as the \emph{base}
and the $\emph{index}$ subexpressions, respectively.
\end{definition}

\ExampleDef{Array Access Expressions}
\listingref{typing-tarray} shows examples of well-typed
array access expressions on the right-hand-side of the
assignments to \verb|int_arr| and \verb|big_little_arr|
and the types inferred for them via the added \verb|as <inferred-type>|.

\RenderProseAndFormally{annotate_expr_EGetArray}
\CodeSubsection{\EGetArrayBegin}{\EGetArrayEnd}{../Typing.ml}

\TypingRuleDef{AnnotateGetArray}
\RenderRelation{annotate_get_array}
See \ExampleRef{Array Access Expressions}.

\RenderProseAndFormally{annotate_get_array}
\CodeSubsection{\AnnotateGetArrayBegin}{\AnnotateGetArrayEnd}{../Typing.ml}

\SemanticsRuleDef{EGetArray}
\ExampleDef{Evaluation of Array Reading Expressions}
In \listingref{semantics-egetarray},
the expression \verb|my_array[[2]]| appearing in the assertion evaluates to the value \texttt{42} since the element
indexed by \texttt{2} in \texttt{my\_array} is \texttt{42}.
\ASLListing{Evaluating an array access expression}{semantics-egetarray}{\semanticstests/SemanticsRule.EGetArray.asl}

\ExampleDef{Evaluation of an Illegal Array Read}
In \listingref{semantics-egetarrayerror},
evaluating the array access expression \verb|my_array[[3]]| results in a \DynamicErrorConfigurationTerm{},
since we are trying to access index \texttt{3} of an array
which has indexes \texttt{0}, \texttt{1} and \texttt{2} only.
\ASLListing{Evaluating an illegal array access}{semantics-egetarrayerror}{\semanticstests/SemanticsRule.EGetArrayTooSmall.asl}

\RenderProseAndFormally{eval_expr_EGetArray}
\CodeSubsection{\EvalEGetArrayBegin}{\EvalEGetArrayEnd}{../Interpreter.ml}

\hypertarget{def-getfieldexpressionterm}{}
\section{Field Reading Expressions\label{sec:FieldReadingExpressions}}
\subsection{Syntax}
\begin{flalign*}
\Nexpr \derives\  & \Nexpr \parsesep \Tdot \parsesep \Tidentifier&
\end{flalign*}

\subsection{Abstract Syntax}
\RenderTypes[remove_hypertargets]{expr_getfield}

\ASTRuleDef{EGetField}
\begin{mathpar}
  \inferrule{
    \buildexpr(\ve) \astarrow \astversion{\ve} \OrBuildError
  }{
  \buildexpr(\overname{\Nexpr(\ve : \Nexpr, \Tdot, \Tidentifier(\id))}{\vparsednode}) \astarrow
  \overname{\EGetField(\astversion{\ve}, \id)}{\vastnode}
}
\end{mathpar}

\subsection{Typing}
\TypingRuleDef{EGetRecordField}
\ExampleDef{Typing Record Field Expressions}
\listingref{eget-record-field} shows field reading expressions
as the \rhsexpressions{} of assignments
and the types inferred for them via the added \verb|as <inferred type>|.
\ASLListing{Typing record field expressions}{eget-record-field}{\typingtests/TypingRule.EGetRecordField.asl}

\ExampleDef{Ill-typed Record Field Expressions}
\listingref{eget-bad-record-field} shows an ill-typed field expression.
\ASLListing{An ill-typed record field expression}{eget-bad-record-field}{\typingtests/TypingRule.EGetBadRecordField.asl}

\RenderProseAndFormally{annotate_expr_EGetField_record_or_exception}

\CodeSubsection{\EGetRecordFieldBegin}{\EGetRecordFieldEnd}{../Typing.ml}
\CodeSubsection{\EGetBadRecordFieldBegin}{\EGetBadRecordFieldEnd}{../Typing.ml}

\TypingRuleDef{EGetCollectionField}

\ExampleDef{Typing Collection Field Expressions}
All of the collection field expressions in
\listingref{typing-lesetcollectionfields} are well-typed.

\RenderProseAndFormally{annotate_expr_EGetField_collection}
\CodeSubsection{\EGetCollectionFieldBegin}{\EGetCollectionFieldEnd}{../Typing.ml}

\TypingRuleDef{EGetBitField}
\ExampleDef{Well-typed Bitfield Expressions}
\listingref{eget-bitfield} shows well-typed bitfield expressions
as the \rhsexpressions{} of assignment statements,
and the types inferred for them via the added \verb|as <inferred-type>|.
\ASLListing{Well-typed bitfield expressions}{eget-bitfield}{\typingtests/TypingRule.EGetBitfield.asl}

\ExampleDef{Nested Bitfield Expressions}
\listingref{eget-bitfield} shows the expression \verb|p.detailed_data.info|,
which refers to the bitfield \verb|info|, which is nested in the bitfield
\verb|detailed_data|.

\ExampleDef{Typed Bitfield Expressions}
\listingref{eget-bitfield} shows the expression \verb|p.detailed_data.info|,
which includes the type annotation \verb|bits(4)|.

\ExampleDef{Ill-typed Bitfield Expressions}
\listingref{eget-bad-bitfield} shows an ill-typed bitfield expression.
\ASLListing{An ill-typed bitfield expression}{eget-bad-bitfield}{\typingtests/TypingRule.EGetBadBitField.asl}

\RenderProseAndFormally{annotate_expr_EGetField_bitfield}
\CodeSubsection{\EGetBitFieldBegin}{\EGetBitFieldEnd}{../Typing.ml}
\CodeSubsection{\EGetBitFieldNestedBegin}{\EGetBitFieldNestedEnd}{../Typing.ml}
\CodeSubsection{\EGetBitFieldTypedBegin}{\EGetBitFieldTypedEnd}{../Typing.ml}
\CodeSubsection{\EGetBadBitFieldBegin}{\EGetBadBitFieldEnd}{../Typing.ml}

\TypingRuleDef{EGetTupleItem}
\ExampleDef{Typing of a Tuple Item Expression}
In \listingref{semantics-egetitem}, the type of the expression \verb|t.item0|
is the \integertypeterm.

In the following rule definition, we use $\itemprefix$ to stand
for its verbatim string.

\RenderProseAndFormally{annotate_expr_EGetField_tuple_item}
\CodeSubsection{\EGetTupleItemBegin}{\EGetTupleItemEnd}{../Typing.ml}

\TypingRuleDef{EGetBadField}
\listingref{eget-bad-bitfield} shows an example of an ill-typed bitfield expression.

\RenderProseAndFormally{annotate_expr_EGetField_error}
\CodeSubsection{\EGetBadFieldBegin}{\EGetBadFieldEnd}{../Typing.ml}

\subsection{Semantics}
\SemanticsRuleDef{EGetField}
\ExampleDef{Evaluation of a Field Read Expression}
In \listingref{semantics-egetfield},
the expression \verb|my_record.a| evaluates to the value \texttt{3}.
\ASLListing{Evaluating a field access expression}{semantics-egetfield}{\semanticstests/SemanticsRule.ERecord.asl}

\RenderProseAndFormally{eval_expr_EGetField}
\CodeSubsection{\EvalEGetFieldBegin}{\EvalEGetFieldEnd}{../Interpreter.ml}

\SemanticsRuleDef{EGetItem}
\ExampleDef{Evaluation of a Tuple Item Expression}
In \listingref{semantics-egetitem},
the expression \verb|t.item0| evaluates to the value \texttt{1}
and
the expression \verb|t.item1| evaluates to the value \texttt{2}.
\ASLListing{Evaluating a tuple item access expression}{semantics-egetitem}{\semanticstests/SemanticsRule.EGetItem.asl}

\RenderProseAndFormally{eval_expr_EGetTupleItem}
\CodeSubsection{\EvalEGetTupleItemBegin}{\EvalEGetTupleItemEnd}{../Interpreter.ml}

\hypertarget{def-getfieldsexpressionterm}{}
\section{Multi-field Reading Expressions\label{sec:MultiFieldReadingExpressions}}
\subsection{Syntax}
\begin{flalign*}
\Nexpr \derives\  & \Nexpr \parsesep \Tdot \parsesep \Tlbracket \parsesep \ClistOne{\Tidentifier} \parsesep \Trbracket &
\end{flalign*}

\subsection{Abstract Syntax}
\RenderTypes[remove_hypertargets]{expr_getfields}


\ASTRuleDef{EGetFields}
\begin{mathpar}
  \inferrule{
    \buildclist[\buildidentity](\vids) \astarrow \vidasts\\
    \buildexpr(\ve) \astarrow \astversion{\ve} \OrBuildError
  }{
    {
      \begin{array}{r}
  \buildexpr(\overname{\Nexpr(\ve : \Nexpr, \Tdot, \Tlbracket, \namednode{\vids}{\ClistOne{\Tidentifier}}, \Trbracket)}{\vparsednode}) \astarrow\\
  \overname{\EGetFields(\astversion{\ve}, \vidasts)}{\vastnode}
      \end{array}
    }
}
\end{mathpar}

\subsection{Typing}
\TypingRuleDef{EGetFields}
\ExampleDef{Typing Multi-field Expressions}
\pagebreak
\listingref{egetfields} shows examples of well-typed multi-field expressions.
\ASLListing{Typing multi-field expressions}{egetfields}{\typingtests/TypingRule.EGetFields.asl}


\RenderProseAndFormally{annotate_expr_EGetFields}
\CodeSubsection{\EGetFieldsBegin}{\EGetFieldsEnd}{../Typing.ml}

\TypingRuleDef{FindBitFieldsSlices}
% Transliteration note: The implementation returns an optional, but all of its uses raise a \typingerrorterm{} if the result is None.
\RenderRelation{find_bitfields_slices}

\ExampleDef{Finding the Slices of a Bitfield}
In \listingref{eget-bitfield},
given the list of bitfields of the \verb|Packet| type,
the list of slices associated with
\verb|data| is \verb|7:1|.
Attempting to obtain the list of slices associated with
\verb|crc|, given the list of bitfields of the \verb|Packet| type,
yields a \typingerrorterm, since it is a nested bitfield.

\RenderProseAndFormally{find_bitfields_slices}

\TypingRuleDef{GetBitfieldWidth}
\RenderRelation{get_bitfield_width}

\ExampleDef{Obtaining the Width of a Bitfield}
In \listingref{egetfields},
obtaining the width of the field \verb|bits3_2|,
given the list of fields of the type \verb|RecordWithBits|
yields the literal expression for \verb|2|.

\RenderProseAndFormally{get_bitfield_width}

\TypingRuleDef{WidthPlus}
% NOTE: the implementation handles just two expressions, here we fold over a list.
\RenderRelation{width_plus}

\ExampleDef{Summing Widths}
Summing the list of expressions \verb|2, 3|,
yields the expression \verb|2+3|,
(in AST terms, \\
$\AbbrevEBinop{\ADD}{\ELInt{2}}{\ELInt{2}}$).

\RenderProseAndFormally{width_plus}

\subsection{Semantics}
\SemanticsRuleDef{EGetfields}
\ExampleDef{Evaluating Multi-field Expressions}
In \listingref{egetfields},
evaluating the multi-field expression \\
\verb|bits_var.[bits3_2, bit1, bit0, info_and_bits]|
yields the bitvector value \\
\verb|'00101010'|,
and evaluating the multi-field expression \\
\verb|record_var.[bits3_2, bit1, bit0, info]|
yields the bitvector value \verb|'001010'|.

\RenderProseAndFormally{eval_expr_EGetFields}
\CodeSubsection{\EvalEGetFieldsBegin}{\EvalEGetFieldsEnd}{../interpreter.ml}

\SemanticsRuleDef{EGetCollectionFields}

\ExampleDef{Typing Collection Fields Expressions}
All of the collection field expressions in
\listingref{typing-lesetcollectionfields} are well-typed.

\RenderProseAndFormally{eval_expr_EGetCollectionFields}
\CodeSubsection{\EvalEGetCollectionFieldsBegin}{\EvalEGetCollectionFieldsEnd}{../interpreter.ml}

\hypertarget{def-atceexpressionterm}{}
\section{Asserting Type Conversion Expressions\label{sec:AssertingTypeConversionExpressions}}

The rule about domains in the definitions of subtype-satisfaction and
type-satisfaction means that it is illegal to use the unconstrained integer
where a constrained integer is expected.
An asserting type conversion, ATC for short, can be used to overcome this.


An ATC allows code to explicitly mark places where uses of constrained types
would otherwise be a static typechecking error. The intent is to reduce the
incidence of unintended errors by making such uses fail typechecking unless
the asserting type conversion is provided.


Note that ATCs are execution-time checks. An execution-time check is a
condition that is evaluated during the evaluation of an execution-time
initializer expression or subprogram. If the condition evaluates to $\False$ it
is a \DynamicErrorConfigurationTerm{}.

\ASLListing{ATC expressions}{typing-atc}{\typingtests/TypingRule.ATC.asl}

\subsection{Syntax}
\begin{flalign*}
\Nexpr \derives\  & \Nexpr \parsesep \Tas \parsesep \Nty &\\
                    |\  & \Nexpr \parsesep \Tas \parsesep \Nconstraintkind &
\end{flalign*}

\subsection{Abstract Syntax}
\RenderTypes[remove_hypertargets]{expr_atc}


\ASTRuleDef{ATC}
\begin{mathpar}
\inferrule[type]{
  \buildexpr(\ve) \astarrow \astversion{\ve} \OrBuildError\hva\\\\
  \buildty(\vt) \astarrow \astversion{\vt} \OrBuildError
}{
  \buildexpr(\overname{\Nexpr(\ve : \Nexpr, \Tas, \vt : \Nty)}{\vparsednode}) \astarrow
  \overname{\EATC(\astversion{\ve}, \astversion{\vt})}{\vastnode}
}
\end{mathpar}

\begin{mathpar}
\inferrule[int\_constraints]{
  \buildexpr(\ve) \astarrow \astversion{\ve} \OrBuildError\hva\\\\
  \buildconstraintkind(\vics) \astarrow \astversion{\vics} \OrBuildError
}{
  {
    \begin{array}{r}
      \buildexpr(\overname{\Nexpr(\ve : \Nexpr, \Tas, \vics : \Nconstraintkind)}{\vparsednode}) \astarrow\\
      \overname{\EATC(\astversion{\ve}, \TInt(\astversion{\vics}))}{\vastnode}
    \end{array}
  }
}
\end{mathpar}

\subsection{Typing}
\TypingRuleDef{ATC}
\ExampleDef{Well-typed Asserting Type Conversion Expressions}
\listingref{typing-atc} shows examples of ATC expressions
and the corresponding annotated expressions in comments.

\ExampleDef{Asserting Type Conversion of Well-constrained Integers}
The specification in \listingref{typing-atc2} is well-typed,
showing how to convert one \\
\wellconstrainedintegertypeterm{} into another
\wellconstrainedintegertypeterm{}.
\ASLListing{ATC of well-constrained integers}{typing-atc2}{\typingtests/TypingRule.ATC2.asl}

\ExampleDef{Safely Applying Asserting Type Conversion to Bitvectors}
The specification in \listingref{typing-atc3} is well-typed,
showing how a \bitvectortypeterm{} may be safely converted
into another \bitvectortypeterm.
\ASLListing{Safely applying ATC to bitvector types}{typing-atc3}{\typingtests/TypingRule.ATC3.asl}

Similarly, the specification in \listingref{typing-atc4} is well-typed,
showing that even literals may require an asserting type conversion to ensure they meet
type-satisfaction requirements.
\ASLListing{Safely applying ATC to a bitvector literal}{typing-atc4}{\typingtests/TypingRule.ATC4.asl}

\ExampleDef{Ill-typed Asserting Type Conversion}
The specification in \listingref{typing-atc-bad} is ill-typed,
as it is applied to unrelated kinds of types.
\ASLListing{Ill-typed ATC of a bitvector type}{typing-atc-bad}{\typingtests/TypingRule.ATC.bad.asl}

\RenderProseAndFormally{annotate_expr_EATC}
\CodeSubsection{\ATCBegin}{\ATCEnd}{../Typing.ml}

\TypingRuleDef{CheckATC}
\texthypertarget{def-compatibletypes}
Intuitively, two types are \emph{compatible} if they are equal up to any constraints.
The formal definition is given next.

\RenderRelation{check_atc}

\ExampleDef{Ill-typed ATC Expressions}
\listingref{checkatc} shows examples of ill-typed ATC expressions.
\ASLListing{Ill-typed ATC expressions}{checkatc}{\typingtests/TypingRule.CheckATC.asl}

\RenderProseAndFormally{check_atc}

\subsection{Semantics}
\SemanticsRuleDef{ATC}
\ExampleDef{Evaluation of an Asserting Type Conversion Expressions}
In \listingref{semantics-atcvalue}, both type assertions ---
\verb|3 as integer| and \verb|3 as integer{3..5}| --- succeed.
\ASLListing{Evaluating a successful type assertion expression}{semantics-atcvalue}{\semanticstests/SemanticsRule.ATCValue.asl}

\ExampleDef{An Unevaluated Asserting Type Conversion}

In \listingref{semantics-atcnotevaluated}, the asserting type conversion on \texttt{y}
does not yield a \DynamicErrorConfigurationTerm{}, since the invocation of \texttt{f1} returns $\False$ when
evaluated:
\ASLListing{An asserting type conversion that is never evaluated}{semantics-atcnotevaluated}
{\semanticstests/SemanticsRule.ATCNotDynamicErrorIfFalse.asl}

\ExampleDef{Asserting Type Conversions with Dynamic Errors}
\listingref{semantics-variousatcs} shows various \dynamicerrorsterm{}.
Note that the \dynamicerrorterm{} appearing in the statement
\verb|var e: integer{4, 5, 6} = 2 as integer{4, 5, 6};|, \\
which is nested in the \verb|if FALSE then|
will never occur, but it is still classified as a \dynamicerrorterm{}.
\ASLListing{Various dynamic errors}{semantics-variousatcs}{\semanticstests/SemanticsRule.ATCVariousErrors.asl}

\RenderProseAndFormally{eval_expr_EATC}
\CodeSubsection{\EvalATCBegin}{\EvalATCEnd}{../Interpreter.ml}

\SemanticsRuleDef{IsValOfType}
\RenderRelation{is_val_of_type}

This relation is used in the context of an asserted type conversion,
which means the typechecker rule \TypingRuleRef{ATC} was already applied,
thus filtering cases where the type inferred for the converted expression
does not type-satisfy $\vt$. The semantics takes this into account and
only returns \False\ in cases where dynamic information is required.

Recall that the $\vt$ is the result of $\annotatetype$, which ensures that
all \\
sub-expressions appearing in $\vt$ are \sideeffectfreeterm{}.

\ExampleDef{Checking Whether a Value Belongs to a Type}
In \listingref{semantics-variousatcs},
checking whether the value \verb|2| is a member of the type \verb|integer{1,2,3}|
succeeds whereas checking whether the value \verb|2| is a member
of the type \verb|integer{4, 5, 6}| fails.

\RenderProseAndFormally{is_val_of_type}
\CodeSubsection{\EvalValOfTypeBegin}{\EvalValOfTypeEnd}{../Interpreter.ml}

\subsubsection{Comments}
Notice that these rules cover all types, including named types ($\TNamed$),
since the \typedast\ returned from \TypingRuleRef{ATC} is the \structureterm\ of the type
given in the specification.
%
Parameterized integers (integers with an empty set of constraints)
cannot appear as a type, since ASL syntax does not allow the following:
\begin{itemize}
\item Declaring an parameterized integer as a variable,
\item Declaring an alias to an parameterized integer type, and
\item Declaring an parameterized integer in a compound type.
\end{itemize}

\SemanticsRuleDef{IsConstraintSat}
\RenderRelation{is_constraint_sat}

See \ExampleRef{Checking Whether a Value Belongs to a Type}.

The use of $\evalexprsef$ in the definition of this function is justified by checks in $\annotatetype$, verifying
that expressions appearing in types are all \sideeffectfreeterm{}.

\RenderProseAndFormally{is_constraint_sat}

\hypertarget{def-patternexpressionterm}{}
\section{Pattern Matching Expressions\label{sec:PatternMatchingExpressions}}
The binary operator $\Tin$ tests whether a value (referred to as the discriminant) matches any item from a $\Npatternset$.
Patterns can also be used to test whether an expression matches a bitmask.
Lists of patterns are also used in case statements.
%
\chapref{PatternMatching} goes into the details of the various types of patterns that can be matched against.

\SyntacticSugarDef{MaskEquality}
The syntax \texttt{x == $m$} is \syntacticsugar{} for \texttt{x IN \{$m$\}},
and it is \desugared{} by \ASTRuleCaseRef{EPattern}{eq}.
The syntax \texttt{x != $m$} is \syntacticsugar{} for \texttt{!x IN \{$m$\}},
and it is \desugared{} by \ASTRuleCaseRef{EPattern}{neq}.

\ASLListing{Pattern expressions}{semantics-epattern}{\semanticstests/SemanticsRule.EPattern.asl}

\subsection{Syntax}
\begin{flalign*}
\Nexpr \derives\  & \Nexpr \parsesep \Tin \parsesep \Npatternset &\\
              |\  & \Nexpr \parsesep \Teqop \parsesep \Tmasklit &\\
              |\  & \Nexpr \parsesep \Tneq \parsesep \Tmasklit &
\end{flalign*}

\subsection{Abstract Syntax}
\RenderTypes[remove_hypertargets]{expr_pattern}


\ASTRuleDef{EPattern}
\begin{mathpar}
\inferrule{
  \buildexpr(\ve) \astarrow \astversion{\ve} \OrBuildError\hva\\\\
  \buildpatternset(\vps) \astarrow \astversion{\vps} \OrBuildError
}{
  {
    \begin{array}{r}
      \buildexpr(\overname{\Nexpr(\ve : \Nexpr, \Tin, \vps : \Npatternset)}{\vparsednode}) \astarrow\\
      \overname{\EPattern(\astversion{\ve}, \astversion{\vps})}{\vastnode}
    \end{array}
  }
}
\end{mathpar}

\begin{mathpar}
\inferrule[eq]{}{
  {
    \begin{array}{r}
  \buildexpr(\overname{\Nexpr(\punnode{\Nexpr}, \Teqop, \Tmasklit(\vm))}{\vparsednode}) \astarrow\\
  \overname{\EPattern(\astof{\vexpr}, \PatternMask(\vm))}{\vastnode}
    \end{array}
  }
}
\end{mathpar}

\begin{mathpar}
\inferrule[neq]{}{
  {
    \begin{array}{r}
      \buildexpr(\overname{\Nexpr(\punnode{\Nexpr}, \Tneq, \Tmasklit(\vm))}{\vparsednode}) \astarrow\\
      \overname{\EPattern(\astof{\vexpr}, \PatternNot(\PatternMask(\vm)))}{\vastnode}
    \end{array}
  }
}
\end{mathpar}

\subsection{Typing}
\TypingRuleDef{EPattern}
\listingref{semantics-epattern} shows examples of pattern expressions.

\RenderProseAndFormally{annotate_expr_EPattern}
\CodeSubsection{\EPatternBegin}{\EPatternEnd}{../Typing.ml}

\subsection{Semantics}
\SemanticsRuleDef{EPattern}
\ExampleDef{Evaluation of Pattern Expressions}
In \listingref{semantics-epattern},
the expression \texttt{42 IN \{0..3, 42\}} evaluates to $\nvbool(\True)$
whereas the expression \texttt{42 IN \{0..3, -4\}} evaluates to $\nvbool(\False)$.
%
In addition, all assertions, which demonstrate pattern expression involving bitmasks,
succeed.

\RenderProseAndFormally{eval_expr_EPattern}
\CodeSubsection{\EvalEPatternBegin}{\EvalEPatternEnd}{../Interpreter.ml}

\hypertarget{def-arbitraryexpressionterm}{}
\section{Arbitrary Value Expressions\label{sec:ArbitraryValueExpressions}}
An expression of the form \texttt{ARBITRARY: ty} evaluates to an arbitrary value in the
domain of \texttt{ty}.
Each evaluation can produce a different arbitrary value, but (as always) once a particular expression is evaluated, its arbitrary value cannot change.
This is because evaluation produces native values, and \ARBITRARY{} is not a valid native value---so once evaluated, it becomes an unchanging native value like any other.

Note that there are two important consequences of producing an arbitrary value when evaluating expressions of the form \texttt{ARBITRARY: ty}:
\begin{enumerate}
  \item The arbitrary value depends only on \texttt{ty}, and no other ASL storage elements.
  \item The only guarantee of the resulting value is that it is a valid member of \texttt{ty}.
    In particular, the language does not define which valid member it is, and ASL specifications must not rely on the value (for example, there is no way to test whether a value was produced by evaluating \ARBITRARY{}).
\end{enumerate}

\subsection{Syntax}
\begin{flalign*}
\Nexpr \derives\  & \Tarbitrary \parsesep \Tcolon \parsesep \Nty &
\end{flalign*}

\subsection{Abstract Syntax}
\RenderTypes[remove_hypertargets]{expr_arbitrary}


\ASTRuleDef{EArbitrary}
\begin{mathpar}
\inferrule{
  \buildty(\vt) \astarrow \astversion{\vt} \OrBuildError
}{
  \buildexpr(\overname{\Nexpr(\Tarbitrary, \Tcolon, \vt : \Nty)}{\vparsednode}) \astarrow
  \overname{\EArbitrary(\astversion{\vt})}{\vastnode}
}
\end{mathpar}

\subsection{Typing}
\TypingRuleDef{EArbitrary}
\ExampleDef{Well-typed Arbitrary Value Expressions}
\listingref{typing-arbitrary} show examples of well-typed arbitrary value expressions.
\ASLListing{Well-typed arbitrary value expressions}{typing-arbitrary}{\typingtests/TypingRule.EArbitrary.asl}

\RenderProseAndFormally{annotate_expr_EArbitrary}
\CodeSubsection{\EArbitraryBegin}{\EArbitraryEnd}{../Typing.ml}

\subsection{Semantics}
\SemanticsRuleDef{EArbitrary}
\ExampleDef{Evaluation of ARBITRARY for an Unconstrained Integer Type}
In \listingref{semantics-earbitraryunconstrained},
the expression \verb|ARBITRARY : integer| evaluates to an arbitrary integer value.
\ASLListing{Evaluating an \texttt{ARBITRARY} expression for an unconstrained integer}{semantics-earbitraryunconstrained}
{\semanticstests/SemanticsRule.EArbitraryInteger3.asl}

\subsubsection{Evaluation of ARBITRARY for a Constrained Integer Type}
In \listingref{semantics-earbitraryconstrained},
the expression \verb|ARBITRARY : integer {3, 42}| evaluates to either \\
$\nvint(3)$ or $\nvint(42)$.
\ASLListing{Evaluating an \texttt{ARBITRARY} expression for a constrained integer}{semantics-earbitraryconstrained}
{\semanticstests/SemanticsRule.EArbitraryIntegerRange3-42-3.asl}

\subsubsection{Evaluation of ARBITRARY for an Integer-indexed Array}
\listingref{semantics-earbitraryarray}
demonstrates how to obtain an arbitrary integer-indexed array, \texttt{int\_array},
and how to obtain an arbitrary enumeration-indexed array, \texttt{enum\_array}.
\ASLListing{Evaluating an \texttt{ARBITRARY} expression for array types}{semantics-earbitraryarray}
{\semanticstests/SemanticsRule.EArbitraryArray.asl}

\RenderProseAndFormally{eval_expr_EArbitrary}
\CodeSubsection{\EvalEArbitraryBegin}{\EvalEArbitraryEnd}{../Interpreter.ml}

\subsubsection{Comments}
Notice that this rule introduces non-determinism.


\section{Structured Type Construction Expressions\label{sec:StructuredTypeConstructionExpressions}}
\hypertarget{def-recordexpressionterm}{}
\listingref{semantics-egetfield} shows an example of a well-typed record construction expression.

\subsection{Syntax}
\begin{flalign*}
\Nexpr \derives\  & \Tidentifier \parsesep \Tlbrace \parsesep \Tminus \parsesep \Trbrace &\\
	       |\ & \Tidentifier \parsesep \Tlbrace \parsesep \ClistOne{\Nfieldassign} \parsesep \Trbrace &\\
\Nfieldassign \derives \ & \Tidentifier \parsesep \Teq \parsesep \Nexpr &
\end{flalign*}

\subsection{Abstract Syntax}
\RenderTypes[remove_hypertargets]{expr_record}


\ASTRuleDef{ERecord}
\begin{mathpar}
  \inferrule[empty]{}{
    {
      \begin{array}{r}
    \buildexpr(\overname{\Nexpr(
    \begin{array}{l}
    \Tidentifier(\vt), \Tlbrace, \Tminus, \Trbrace
    \end{array}
    )}{\vparsednode}) \\
    \astarrow\ \overname{\ERecord(\TNamed(\vt), \emptylist)}{\vastnode}
\end{array}
}
}
\end{mathpar}

\begin{mathpar}
  \inferrule[non\_empty]{
    \buildclist[\buildfieldassign](\vfieldassigns) \astarrow \vfieldassignasts
  }{
    {
      \begin{array}{r}
  \buildexpr\left(\overname{\Nexpr\left(
    \begin{array}{l}
    \Tidentifier(\vt), \Tlbrace, \\
    \wrappedline\ \namednode{\vfieldassigns}{\ClistOne{\Nfieldassign}}, \\
    \wrappedline\ \Trbrace
    \end{array}
    \right)}{\vparsednode}\right) \\
    \astarrow\ \overname{\ERecord(\TNamed(\vt), \vfieldassignasts)}{\vastnode}
\end{array}
}
}
\end{mathpar}

\ASTRuleDef{FieldAssign}
\NotImportedToASLSpecYet{
\hypertarget{build-fieldassign}{}
The function
\[
  \buildfieldassign(\overname{\parsenode{\Nfieldassign}}{\vparsednode}) \;\aslto\; \overname{(\Identifier, \expr)}{\vastnode}
\]
transforms a parse node $\vparsednode$ into an AST node $\vastnode$.
} % END_OF_UNIMPORTED_RELATION


\begin{mathpar}
\inferrule{}{
  \buildfieldassign(\Nfieldassign(\Tidentifier(\id), \Teq, \punnode{\Nexpr})) \astarrow
  \overname{(\id, \astof{\vexpr})}{\vastnode}
}
\end{mathpar}

\subsection{Typing}
\TypingRuleDef{ERecord}
\listingref{semantics-egetfield} shows an example of a well-typed record construction expression
and an example (in comment) where the same field is initialized twice, which is invalid.

\RenderProseAndFormally{annotate_expr_ERecord}
\CodeSubsection{\ERecordBegin}{\ERecordEnd}{../Typing.ml}


\TypingRuleDef{AnnotateFieldInit}
\RenderRelation{annotate_field_init}

See \listingref{semantics-egetfield} for an example.

\RenderProseAndFormally{annotate_field_init}

\subsection{Semantics}
\SemanticsRuleDef{ERecord}
\ExampleDef{Evaluation of Record Construction Expressions}
In \listingref{semantics-egetfield},
the expression \verb|MyRecordType{a=3, b=42}| evaluates to the native record value \\
$\NVRecord(\va\mapsto\nvint(3), \vb\mapsto\nvint(42))$.

\RenderProseAndFormally{eval_expr_ERecord}
\CodeSubsection{\EvalERecordBegin}{\EvalERecordEnd}{../Interpreter.ml}

\section{Tuple Expressions\label{sec:TupleExpressions}}
\hypertarget{def-tupleexpressionterm}{}

\ASLListing{Tuple expressions}{TupleExpressions}{\definitiontests/TupleExpressions.asl}

\subsection{Syntax}
\begin{flalign*}
\Nexpr \derives\  & \Plisttwo{\Nexpr} &
\end{flalign*}

\subsection{Abstract Syntax}
\RenderTypes[remove_hypertargets]{expr_tuple}


\ASTRuleDef{ETuple}
\begin{mathpar}
\inferrule{
  \buildplist[\buildexpr](\vexprs) \astarrow \vexprasts
}{
  \buildexpr(\overname{\Nexpr(\namednode{\vexprs}{\Plisttwo{\Nexpr}})}{\vparsednode}) \astarrow
  \overname{\ETuple(\vexprasts)}{\vastnode}
}
\end{mathpar}

\subsection{Typing}
\TypingRuleDef{ETuple}
\listingref{semantics-etuple} shows an example of a well-typed tuple expressions.

\RenderProseAndFormally{annotate_expr_ETuple}
\CodeSubsection{\ETupleBegin}{\ETupleEnd}{../Typing.ml}

\subsection{Semantics}
\SemanticsRuleDef{ETuple}
\ExampleDef{Evaluation of Tuple Expressions}
In \listingref{semantics-etuple},
the expression \texttt{(3, Return42())} evaluates to the value \texttt{(3, 42)}.

\ASLListing{Evaluation of tuple expressions}{semantics-etuple}{\semanticstests/SemanticsRule.ETuple.asl}

\RenderProseAndFormally{eval_expr_ETuple}
\CodeSubsection{\EvalETupleBegin}{\EvalETupleEnd}{../Interpreter.ml}

\section{Parenthesized Expressions\label{sec:ParenthesizedExpressions}}
A single expression inside parentheses is not considered to be a tuple, but rather the element
inside the parenthesis.
Parenthesizing an expression can be used to improve readability, enforce an order of evaluation,
and avoid binary operator precedence errors (see \ASTRuleRef{CheckNotSamePrec}).

\subsection{Syntax}
\begin{flalign*}
\Nexpr \derives\  & \Tlpar \parsesep \Nexpr \parsesep \Trpar &
\end{flalign*}

\subsection{Abstract Syntax}
We represent a parenthesized expression as a single element tuple for the technical
reason of enabling \ASTRuleRef{CheckNotSamePrec} by distinguishing parenthesized
expressions from non-parenthesized expressions. However, upon typing such expressions,
we ignore the parenthesis (see \TypingRuleRef{ETuple}.\textsc{PARENTHESIZED}).

\ASTRuleDef{ParenExpr}
\begin{mathpar}
  \inferrule[sub\_expr]{}{
  \buildexpr(\overname{\Nexpr(\Tlpar, \punnode{\Nexpr}, \Trpar)}{\vparsednode}) \astarrow
  \overname{\ETuple([\ \astof{\vexpr}\ ])}{\vastnode}
}
\end{mathpar}

\section{Array Construction Expressions\label{sec:ArrayConstructionExpressions}}
Array construction expression are used by the type system to express the initialization
of array-typed variables. Since there is no syntax to initialize arrays, there are also
no rules for building the AST for such expressions nor rules for typechecking them.

\subsection{Abstract Syntax}
\RenderTypes[remove_hypertargets]{expr_array}

\subsection{Semantics}

The Semantic Rules use $\evalexprsef\empty$ because the typechecker in
$\annotatetype$ guarantees that expressions in types are \sideeffectfreeterm{}.

\SemanticsRuleDef{EArray}
In \listingref{typing-tarray},
the variable \verb|int_arr| is initialized with an array construction expression
(which is not expressible in ASL text, only in \typedast, but conceptually could be
represented as \verb|array[[4]] of 0|),
which evaluates to
\[
\NVVector([\nvint(0), \nvint(0), \nvint(0), \nvint(0)]) \enspace.
\]

\RenderProseAndFormally{eval_expr_EArray}
\CodeSubsection{\EvalEArrayBegin}{\EvalEArrayEnd}{../Interpreter.ml}

\section{Side-effect-free Expressions\label{sec:SideEffectFreeExpressions}}
\subsection{Typing}
\texthypertarget{def-sideeffectfreeexpressionterm}
An expression $\ve$ is considered to be \sideeffectfreeterm{} in the \staticenvironmentterm{} $\tenv$
if $\annotateexpr(\tenv, \ve) \typearrow (\Ignore, \Ignore, \vses)$
and $\vses$ is \readonlyterm{}.

\subsection{Semantics}
\SemanticsRuleDef{ESideEffectFreeExpr}
\RenderRelation{eval_expr_sef}

\ExampleDef{Evaluating a Side-effect-free Expression}
Evaluating the literal expression $\LInt(1)$ in any environment $\env$,
yields \\
$\ResultExprSEF(\nvint(1), \emptygraph)$.

\RenderProseAndFormally{eval_expr_sef}
Notice that the output configuration does not contain an environment,
since \sideeffectfreeexpressionsterm{} do not modify the environment.

\section{Evaluating a List of Expressions\label{sec:ExprList}}
\SemanticsRuleDef{EExprList}
\RenderRelation{eval_expr_list}

\ExampleDef{Evaluating a List of Expressions}
In \listingref{semantics-egetfield},
evaluating the expression \verb|MyRecordType{a=3, b=42}| entails evaluating
the expression list \verb|3, 42|, which yields
the list of values $[\nvint(3), \nvint(42)]$.

\RenderProseAndFormally{eval_expr_list}
